\documentclass[11pt]{article}
\usepackage[T1]{fontenc}
\usepackage{lmodern}
\usepackage{amsmath,amssymb,amsthm}
\usepackage[margin=1in]{geometry}
\usepackage[hidelinks]{hyperref}
\usepackage{microtype}

\newtheorem{theorem}{Theorem}
\newcommand{\R}{\mathbb{R}}
\newcommand{\Cut}{\operatorname{Cut}}

\title{A Sharp Nonattained Clique--Web Ratio\\
on a Three-Parameter Lorentzian Family}
\author{Xiaojian Zeng\\
PBC School of Finance, Tsinghua University, 100190\\
\texttt{xiaojianzeng@foxmail.com}}
\date{}
\hypersetup{
  pdftitle={A Sharp Nonattained Clique-Web Ratio on a Three-Parameter Lorentzian Family},
  pdfauthor={Xiaojian Zeng},
  pdfsubject={Sharp optimization on a Lorentzian matrix family}
}

\begin{document}
\maketitle

\begin{abstract}
For an explicit three-parameter family of normalized positive Lorentzian
matrices associated with a seven-vertex clique--web exponent, we determine
the ratio supremum exactly.  It is
\(\bigl((5+\sqrt5)/4\bigr)^{5/4}\), is not attained at any finite
parameter triple, and is approached by a closed-form sequence lying on both
boundary equations of the Lorentzian region.  We also recover the published
feasibility criterion directly from an invariant-subspace decomposition.
The result concerns only this family and does not determine the unrestricted
bounded-ratio constant.
\end{abstract}

\section{Context and statement}

Lorentzian polynomials provide a broad setting in which a Hessian has a
single positive direction; see Br\"and\'en and Huh~\cite{BrandenHuh2020}.
For a real symmetric matrix with nonnegative entries, we use the corresponding
quadratic notion: the matrix is \emph{Lorentzian} if it has at most one
positive eigenvalue.  Huang, Huh, Soskin, and Wang~\cite{HuangEtAl2025}
introduced bounded ratios for such matrices, identified their cone with the
dual of the cut cone, and conjectured that every normalized reduced ratio is
bounded by~2.  Clique--web facets belong to the standard structure theory of
cut cones~\cite[Sections~29.3--29.6]{DezaLaurent1997}.

Baldi and Kummer~\cite{BaldiKummer2026} used a seven-vertex clique--web facet
to disprove that conjectured bound.  Their construction uses, for
\(r,q,t>0\),
\[
P(r,q,t)=
\begin{pmatrix}
1&r&q&q&r&t&t\\
r&1&r&q&q&t&t\\
q&r&1&r&q&t&t\\
q&q&r&1&r&t&t\\
r&q&q&r&1&t&t\\
t&t&t&t&t&1&1\\
t&t&t&t&t&1&1
\end{pmatrix}.
\]
Let
\[
 z=x_{13}+x_{14}+x_{24}+x_{25}+x_{35}+x_{67}
   -\sum_{i=1}^{5}(x_{i6}+x_{i7}),
 \qquad \alpha=\frac{z}{4}.
\]
The coordinate sum of \(\alpha\) is \(-1\), and direct substitution gives
\begin{equation}\label{eq:objective}
 P(r,q,t)^\alpha
 =\left(\frac{q^5}{t^{10}}\right)^{1/4}
 =\left(\frac{q}{t^2}\right)^{5/4}.
\end{equation}
Baldi and Kummer proved that this matrix is Lorentzian precisely when
\begin{equation}\label{eq:feasible}
 r+q-2\geq \sqrt5\,|r-q|,
 \qquad
 t^2\geq \frac{1+2r+2q}{5},
\end{equation}
and exhibited the feasible point \((284,701,20)\), whose value exceeds~2.
They did not optimize the displayed family.

Restricting a matrix inequality to a structured normalized family and then
determining its exact sharp ratio is also useful in other settings.  For
example, Zeng, Liu, and Ratnavelu~\cite{ZengEtAl2026} determine a sharp
permanent ratio on rank-two real correlation matrices by reducing to a
low-dimensional exact minimization.  That positive-semidefinite permanent
problem is otherwise unrelated to the Lorentzian ratio considered here.

Define the family constant
\[
 C_{\mathrm{fam}}
 :=\sup\left\{P(r,q,t)^\alpha:
 r,q,t>0\ \text{and \eqref{eq:feasible} holds}\right\}.
\]
Our result resolves exactly this restricted optimization.

\begin{theorem}\label{thm:main}
One has
\[
 C_{\mathrm{fam}}=\left(\frac{5+\sqrt5}{4}\right)^{5/4}.
\]
No finite feasible triple attains this value.  For every integer \(n\geq1\),
set
\begin{equation}\label{eq:sequence}
 q_n=n,\qquad
 r_n=\frac{(\sqrt5-1)+(3-\sqrt5)n}{2},\qquad
 t_n=\sqrt{\frac{\sqrt5+(5-\sqrt5)n}{5}}.
\end{equation}
Then \((r_n,q_n,t_n)\) is feasible and
\[
 P(r_n,q_n,t_n)^\alpha
 =\left(\frac{5n}{\sqrt5+(5-\sqrt5)n}\right)^{5/4}
 \longrightarrow
 \left(\frac{5+\sqrt5}{4}\right)^{5/4}.
\]
\end{theorem}

\section{The Lorentzian region}

For completeness, we independently recover~\eqref{eq:feasible}.  Let
\[
 U=\langle e_6-e_7\rangle,
 \quad
 V=\langle e_6+e_7,e_1+\cdots+e_5\rangle,
 \quad
 W=\left\{(w_1,\ldots,w_5,0,0):\sum_{i=1}^5w_i=0\right\}.
\]
These subspaces give a direct sum of \(\R^7\) and are mutually orthogonal
for the bilinear form represented by \(P\).  The restriction to \(U\) is
zero.  In the displayed basis of \(V\), the Gram matrix is
\[
 \begin{pmatrix}
 4&10t\\
 10t&5+10r+10q
 \end{pmatrix}.
\]
It has a positive direction.  Hence its restriction has at most one positive
eigenvalue exactly when its determinant is nonpositive, or equivalently when
\[
 t^2\geq\frac{1+2r+2q}{5}.
\]

The upper-left \(5\times5\) block is circulant.  On \(W\), its two distinct
eigenvalues, each of multiplicity two, are
\[
 1-\frac{r+q}{2}\mathbin{\pm}\frac{\sqrt5}{2}(r-q).
\]
Both are nonpositive exactly when
\(r+q-2\geq\sqrt5|r-q|\).  Since the restriction to \(V\) always has a
positive direction, the two inequalities in~\eqref{eq:feasible} are together
necessary and sufficient for \(P\) to have exactly one positive eigenvalue.
This argument includes the equality cases and the unavoidable kernel on
\(U\).

\section{Sharp optimization and nonattainment}

\begin{proof}[Proof of Theorem~\ref{thm:main}]
Put \(s=\sqrt5\).  By~\eqref{eq:objective}, for fixed feasible \((r,q)\)
the objective is strictly decreasing in \(t^2\).  Thus it suffices to take
\[
 t^2=\frac{1+2r+2q}{5}
 \quad\text{and maximize}\quad
 B(r,q):=\frac{5q}{1+2r+2q}.
\]

First suppose that \(r\leq q\).  The first inequality
in~\eqref{eq:feasible} is then equivalent to
\[
 r\geq a(q):=\frac{(s-1)q+2}{s+1}
 =\frac{(3-s)q+(s-1)}{2}.
\]
The interval \([a(q),q]\) is nonempty exactly when \(q\geq1\).  Because
\(B(r,q)\) decreases with \(r\),
\begin{equation}\label{eq:firstbranch}
 B(r,q)\leq B(a(q),q)
 =\frac{5q}{s+(5-s)q}
 <\frac{5}{5-s}.
\end{equation}
The last inequality is strict for every finite \(q>0\), since the denominator
on its left contains the additional positive term~\(s\).

Now suppose that \(r\geq q\).  Feasibility is equivalent to
\[
 (s-1)r\leq(s+1)q-2,
\]
and compatibility with \(r\geq q\) again forces \(q\geq1\).  Therefore
\begin{equation}\label{eq:secondbranch}
 B(r,q)\leq\frac{5q}{1+4q}<\frac{5}{5-s}.
\end{equation}
Indeed, the last comparison is equivalent to
\((5-s)q<1+4q\), which is immediate.  The two branches exhaust the feasible
set, so every finite feasible triple satisfies
\[
 P(r,q,t)^\alpha
 <\left(\frac{5}{5-s}\right)^{5/4}
 =\left(\frac{5+s}{4}\right)^{5/4}.
\]
This proves the upper bound and already proves nonattainment.

It remains to show sharpness.  For~\eqref{eq:sequence}, one has
\[
 q_n-r_n=\frac{(s-1)(n-1)}2,
 \qquad
 r_n+q_n-2=\frac{(5-s)(n-1)}2=s(q_n-r_n).
\]
Also \(q_n\geq r_n>0\), \(t_n>0\), and direct substitution gives
\[
 t_n^2=\frac{1+2r_n+2q_n}{5}.
\]
Thus both feasibility inequalities hold with equality.  Finally,
\[
 P(r_n,q_n,t_n)^\alpha
 =\left(\frac{5n}{s+(5-s)n}\right)^{5/4}
 \longrightarrow
 \left(\frac{5}{5-s}\right)^{5/4},
\]
which supplies the matching lower bound.
\end{proof}

\section{Scope and comparison}

The finite point of Baldi and Kummer gives
\((701/400)^{5/4}\approx2.01638\).  Theorem~\ref{thm:main} replaces that
single witness, within the same three-parameter family, by the exact family
supremum, an attainment classification, and an exact extremal sequence.  No
claim is made that \(C_{\mathrm{fam}}\) equals the global bounded-ratio
constant \(f(\alpha)\), nor does this calculation provide a global upper
bound for that constant.

\paragraph{Computational and AI assistance.}
AI-assisted tools were used for literature triage, symbolic cross-checks,
and preparation of the \LaTeX{} manuscript.  Exact algebraic identities and
the mathematical argument were checked independently; the proof above is
self-contained and does not rely on floating-point computation.  The author
takes responsibility for all statements and citations.

\bibliographystyle{plain}
\bibliography{references}

\end{document}
